% !TeX root = ../thesis.tex
% Vorläufige Gliederung mit geplanten Inhalten.
\chapter{Technische Umsetzung des Evaluationsframeworks}
\label{chap:evaluationsframework}
\section{Aufbau und Systemarchitektur}
\begin{inhalte}
  \item Modularer Aufbau
  \item Einheitliche Modell- beziehungsweise Pipeline-Schnittstellen
  \item Reproduzierbare Verarbeitung
  \item Einlesen der Dokumente und Bildaufbereitung
  \item Gegebenenfalls Feld- oder Bereichserkennung
\end{inhalte}
\section{Implementierung der untersuchten Ansätze}
\begin{inhalte}
  \item Integration spezialisierter \acs{ocr}-/\acs{htr}- beziehungsweise Machine-Learning-Ansätze
  \item Integration von Vision-Language-Modellen: Prompting, Bildverarbeitung, strukturierte Extraktion und Ausgabe im Format \ac{json}
  \item Kombination mehrerer Modelle beziehungsweise Verarbeitungsschritte
  \item Beispielsweise \acs{ocr}/\acs{htr} mit anschließender Strukturierung oder Korrektur durch ein \ac{llm}
  \item Standardisierte Ergebnisstruktur
\end{inhalte}
\section{Automatisierte Durchführung und Evaluation}
\begin{inhalte}
  \item Automatisierter Ground-Truth-Vergleich
  \item Speicherung von Modellantworten und Metriken
  \item Erfassung von Laufzeit und Kosten
  \item Gegebenenfalls Kennzeichnung unsicherer Extraktionen anhand einfacher Konsistenz- oder Confidence-Verfahren
\end{inhalte}
